According to the cited press reports, on Thursday 8 October 2026, at about 17:00 (UTC$-$3), a flow of water, mud and rocks descended from the ravine draining the Andean foothills of San Carlos de Apoquindo (Las Condes, Santiago) and traveled along residential streets for several kilometers. It swept away vehicles and a bus \citep{theclinic2026frontal}, flooded homes, a clinic and a university, injured six people and prompted a disaster-area declaration \citep{interferencia2026}. According to the press reports consulted, which cite information from Servicio Nacional de Geología y Minería (Sernageomin), the event occurred in an area classified in February 2026 as having a high debris-flow hazard \citep{chilevision2026}. According to the municipality and BioBioChile, the catchment had been affected by a fire covering about \SI{1000}{ha} between December 2025 and January 2026 \citep{lascondes2026,biobio2026origen}. According to CR2, the event followed an exceptionally rainy start to October: more than \SI{90}{mm} at Quinta Normal so far that month, close to the historical maximum of October 1930 \citep{cr2_2026}.

By the following day, no official mapping of the affected area, discharge measurements or published high-water marks had been located among the public sources consulted. The available information was limited to press reports, the municipality's list of closed streets, an affected-area perimeter and two conflicting source accounts of the origin. This work tests scenarios of the debris-flow path, volume and discharge using an explicit physical model and verifiable public evidence. For those scenarios, it also calculates flow depth, velocity and pressure in streets and at buildings in the study area, taking the uncertainty of the estimates into account.
